\label{combinatorialIntervalSection}
In this section, we give combinatorial results on the chain complex corresponding to an interval. 

\begin{remark}
  Fix $n:\mathbb N$. Consider the simplicial complex $\mathbb I_n^\sim$ whose vertices are elements of $\mathrm{Fin}(2^n)$ and only contains simplices of numbers whose pairwise distance is $\leq 1$. 
  Note that the set $\mathbb I_n^{\sim k}$ in \Cref{description-Cn-simn} is the set of elementary $k$-chains of $\mathbb I_n^\sim$ as in \parencite[Def VI.2.1]{EilenbergSteenrod}.
\end{remark}   

\begin{definition}{\cite[Def VI.2.3]{EilenbergSteenrod}}
  We define the ordered chain complex $C_k$ associated to $\mathbb I_n^{\sim}$ as follows: for each $k$, the group $C_k$ is given by the free abelian group on $\mathbb I_n^{\sim k}$, and the map $\partial_k: C_k\to C_{k-1}$ is induced by 
  $$\partial_k(a_0,\hdots , a_k) = \sum\limits_{i=0}^k (-1)^i(a_0,\hdots, \widehat{a_i},\hdots, a_k).$$

  We define the alternating chain complex $D_k$ associated to $\mathbb I_n^\sim$ as follows: for each $k$, the group $D_k$ is obtained as quotient of $C_k$ by setting
  \begin{align*}
    (a_0, \hdots, a_k) - \mathrm{sgn}(\sigma) (a_{\sigma(0)} , \hdots, a_{\sigma(k)})=0 &\text{ for all permutations } \sigma \text{ on } \mathbb N_{\leq k} \\
    (a_0, \hdots, a_k) = 0 &\text{ whenever } a_i = a_j \text{ for some } i\neq j
  \end{align*}
  where $\mathrm{sgn}(\sigma)$ is the sign of the permutation $\sigma$.
\end{definition}

Using \cite[Lemmas VI.6.5 and VI.6.6]{EilenbergSteenrod}, one can see that the alternating chain complex defined above agrees with the one defined in \cite[Def VI.6.8]{EilenbergSteenrod}.

\begin{remark}\label{RmkTrivialDk}
As any tuple in $\mathbb I_n^{\sim k}$ contains at most two distinct elements, $D_k$ is trivial for $k\geq 2$. 
\end{remark}

\begin{definition}\label{DfnExactSequence}
For a given group $G$ and natural number $n$, we are interested in the following sequence:
\begin{equation}\label{EqSequence}
    \begin{tikzcd}
      0 \arrow[r,"d_{-2}"]& G \arrow[r,"d_{-1}"]& 
      G^{\mathbb I_n^{\sim 0}  } \arrow [r , "d_0"]     & 
      G^{\mathbb I_n^{\sim 1}  } \arrow [r , "d_1"]     & 
      \cdots                     \arrow [r , "d_{k-2}"] & 
      G^{\mathbb I_n^{\sim k-1}} \arrow [r , "d_{k-1}"] & 
      G^{\mathbb I_n^{\sim k  }} \arrow [r , "d_k"]     & 
      G^{\mathbb I_n^{\sim k+1}} \arrow [r , "d_{k+1}"] & 
      \cdots
    \end{tikzcd}
  \end{equation}
  with $d_k(f)(a_0, \hdots , a_k) = \prod\limits_{i=0}^k f(a_0 , \hdots, \widehat {a_i} , \hdots, a_k)^{(-1)^i}$. Note that in general, this sequence is a sequence of pointed sets. If $G$ is an abelian group and we have the usual alternating sum formula for $d_k$. 
\end{definition}

\begin{lemma}\label{LemmaExactnessIntervalApproxFirstTerms}
  For any group $G$, the sequence in \Cref{EqSequence} is exact at 
  $G,G^{\mathbb I_n^{\sim 0}},$ and $G^{\mathbb I_n^{\sim 1}}$.
\end{lemma}

\begin{proof}
  \begin{itemize}
    \item
  As $\mathbb I_n$ is inhabited, $G \to G^{\mathbb I_n^{\sim0}}$ is injective, thus we have exactness at $G$. 
%
\item 
Suppose $\alpha : \mathbb I_n \to G$ is a cocycle. 
  Then $\alpha(a) = \alpha(b)$ whenever $|a - b|\leq 1$. 
%  By \Cref{RmkConstantIfRespsimn}, $\alpha$ is constant.
  But as all elements of $\mathrm{Fin}(2^n)$ are related by steps of size $1$, $\alpha$ is constant.
  Thus we have exactness at $G^{\mathbb I_n^{\sim 0}}$. 
%
\item
  If $\alpha : \mathbb I_n^{\sim 1} \to G$ is a cocycle, then 
  $\alpha(a,a) = 1$ and $\alpha(a,b) = \alpha(b,a)^{-1}$ for all 
  $(a,b): \mathbb I_n^{\sim 1}$. Define 
  $\beta : \mathbb I_n^{\sim 0} \to G$ by 
 $$\beta(a) = \alpha(a-1,a)\cdot \alpha (a-2,a-1) \cdot \hdots \cdot\alpha(0,1).$$
  Then $\beta(a) \cdot \beta(a)^{-1} = 1 = \alpha(a,a)$, 
  and $\beta(a+1) \cdot \beta(a)^{-1} = \alpha(a,a+1)$, and hence 
  $\beta(a)\cdot \beta(a+1)^{-1} = \alpha(a,a+1)^{-1} = \alpha(a+1,a)$. 
  Thus $d_0(\beta) $ and $\alpha$ agree on all of $\mathbb I_n^{\sim 1}$
  and we have exactness at $G^{\mathbb I_n^{\sim1}}$.
\end{itemize}
\end{proof}

 \begin{lemma}\label{LemmaExactnessIntervalApproxLaterTerms}
   For any abelian group $G$, the sequence in \Cref{EqSequence} is exact at $G^{\mathbb I_n^{\sim k}}$
  for $k\geq 2$. 
\end{lemma} 

\begin{proof}
  The sequence in \Cref{EqSequence} is given by applying $\mathrm{Hom}_{\mathrm{Ab}}(-,G)$ to the ordered chain complex $C_k$.
  By \parencite[Theorem VI.6.10]{EilenbergSteenrod}, this ordered chain complex is homotopy equivalent to the alternating chain complex $D_k$. Then $\mathrm{Hom}_{\mathrm{Ab}}(D_k,G)$ is trivial for $k\geq 2$, as $D_k$ was trivial by \Cref{RmkTrivialDk}.
  We conclude that the sequence in \Cref{EqSequence} is exact at $G^{\mathbb I_n^{\sim k}}$.
\end{proof}


